
\section{模型的改进与推广}

\subsection{模型的改进}

本文的几何模型假设金属A为严格的直圆柱体、金属B为严格的球体，在后续改进中可引入更真实的颗粒形状如哑铃形、椭球形或带有表面粗糙度的颗粒，以逼近实际导电填料的复杂形貌。蒙特卡洛模拟方面，当前采用简单的随机均匀采样生成颗粒，当填充率较高时颗粒间的体积排斥效应不容忽视，可引入Metropolis-Hastings算法或分子动力学弛豫步骤确保生成配置的物理合理性。连通性判定方面，当前仅考虑了两颗粒间的二元连通性，可通过引入接触电阻网络模型将二元连通扩展为加权连通的导纳网络，进而计算微构体的宏观等效电导率，而不仅限于导通或不通导的二值判定。计算效率方面，可进一步采用GPU并行化加速颗粒间距离计算，支持更大规模（$N > 10^4$）的高精度模拟。

\subsection{模型的推广}

本文提出的三维几何建模与蒙特卡洛模拟框架具有较强的通用性，可推广至以下场景。一是导电高分子复合材料的设计优化，将碳纳米管类比为金属A（高长径比圆柱体）、将炭黑颗粒类比为金属B（近球形填料），复用本文的渗透阈值和成本优化分析方法，指导工业级导电塑料的配方设计。二是锂离子电池电极材料中导电剂网络的优化设计，导电剂颗粒在活性材料颗粒间形成导电通路的问题与本文的微构体导通问题在几何本质上高度相似。三是多孔介质中的流体渗透问题，将导电通路替换为流体通路，隧穿距离替换为孔隙喉道尺寸，即可将本文框架推广至油气储层岩石的渗透率预测。推广时需根据具体物理场景调整颗粒几何参数、连通判据和边界条件。
